今後の研究・実装に関する方向性を以下に示す。特に、精度向上だけでなく、実運用での信頼性・説明可能性・運用コストを総合的に高めることを重視する。

\begin{enumerate}
\item \textbf{重み付けと学習履歴の動的最適化：}本実験2で示唆された最適学習件数（15〜20件）と時間減衰の効果を踏まえ、ワーカーごとの利用頻度や行動変化に応じて学習窓幅と重みを自動調整する。回避業務ペナルティやリピート志向の強さも含め、ユーザー単位での適応的スコアリングを実装する。
\item \textbf{モデル選定と推論コストの最適化：}実験3の比較検討を拡充し、日本語モデルの精度・推論速度・運用コストを軸に評価基準を整備する。モデルの蒸留や軽量化、事前エンコードの高速化により、リアルタイム推薦の安定性を高める。
\item \textbf{実サービス検証と評価設計：}A/Bテストにより応募率・継続率・定着率の改善を定量化し、推薦のROIを示す。短期KPI（応募・閲覧・クリック）と中期KPI（勤務実績・リピート率）を分離して評価し、最適化の指標を明確化する。
\item \textbf{説明可能性とUXの向上：}推薦理由の可視化（例：業務内容一致、リピート施設、移動コストの低さ）を提示し、ワーカーと事業者双方の納得感を高める。総額志向に合わせた報酬表示や、回避業務の注意喚起などUI側の改善と連動させる。
\item \textbf{データ品質と運用ルールの整備：}求人情報の記述揺れや欠損、施設ごとの入力品質の差が推薦精度に影響するため、入力補助や正規化ルールを整備する。回避業務の更新頻度や、口コミ・評価データの取り扱い指針を含め、運用面でのガバナンスを確立する。
\item \textbf{多面的特徴量の統合：}口コミ、評価、職場環境、教育体制、オンボーディング情報など、非テキスト要因を特徴量化して統合する。これにより短期的な適合だけでなく、中長期の満足度・定着率に寄与する推薦へと拡張する。
\end{enumerate}
